\section*{集合・命題}

\subsection*{この分野で起こりやすいミス}
\begin{itemize}
  \item 条件をそのまま読むだけで，倍数・約数・素因数の関係に言い換えられない。
  \item 必要条件と十分条件の向きを逆にとらえてしまう。
  \item $A\cap \overline{B}$ などの記号を，意味ではなく雰囲気で処理してしまう。
\end{itemize}

\subsection*{例題}
$U=\{1,2,3,\dots,20\}$ とする。  
$a=3,\ b=4$ のとき，
\[
A=\text{$a$ の倍数全体の集合},\qquad B=\text{$b$ を約数にもたない数の集合}
\]
のような設定や，
\[
A=\text{3の倍数の集合},\qquad B=\text{2の倍数の集合}
\]
として，$A\cap B,\ A\cap \overline{B}$ を求める問題。

\subsection*{例題で起こりがちなミス}
\begin{itemize}
  \item $2$でも$3$でも割り切れることから，すぐに「$6$の倍数」とまとめられない。
  \item $A\cap \overline{B}$ を「$A$に属し，$B$には属さないもの」と読めない。
  \item 集合を具体的に書き出さず，頭の中だけで処理して取り違える。
\end{itemize}

\subsection*{解説を読むときに注意してほしいこと}
\begin{itemize}
  \item まず「この集合は何を表しているのか」を日本語で言い換えてから読む。
  \item 解説の中でいきなり集合が書かれていても，「なぜその要素だけが残るのか」を1つずつ確認する。
  \item 記号操作としてではなく，「共通部分」「補集合」の意味を毎回言葉でとらえる。
\end{itemize}

\subsection*{今後の学習で意識するポイント}
\begin{itemize}
  \item 倍数・約数・素因数の言い換えをすぐにできるようにする。
  \item ベン図や要素の書き出しを省略しない。
  \item 条件を見たら，まず具体例を2つか3つ試して意味を確認する習慣をつける。
\end{itemize}

\bigskip

\section*{図形と計量}

\subsection*{この分野で起こりやすいミス}
\begin{itemize}
  \item どの公式を使えばよいか判断できない。
  \item 四角形や複雑な図形を三角形に分ける発想が出てこない。
  \item $\sin(180^\circ-\theta)=\sin\theta$ などの補角の関係を使えない。
\end{itemize}

\subsection*{例題}
四角形 $ABCD$ において，対角線で分けた2つの三角形の面積を
\[
S_1=\triangle ABD,\qquad S_2=\triangle BCD
\]
として，
\[
S_1=\frac12 AB\cdot AD\sin A,\qquad
S_2=\frac12 BC\cdot CD\sin C
\]
を利用し，
\[
A+C=180^\circ
\]
から
\[
\sin C=\sin A
\]
を用いて式を整理する問題。

\subsection*{例題で起こりがちなミス}
\begin{itemize}
  \item 四角形のまま扱おうとして手が止まる。
  \item $\sin C=\sin A$ に変形できず，式の整理が進まない。
  \item 面積公式
  \[
  S=\frac12 ab\sin\theta
  \]
  を使う場面だと気づけない。
\end{itemize}

\subsection*{解説を読むときに注意してほしいこと}
\begin{itemize}
  \item 「なぜ四角形を三角形2つに分けたのか」を確認する。
  \item 面積公式を使っている箇所では，どの2辺とどの角を使っているのかを図と対応させる。
  \item $\sin C$ を $\sin A$ に変えている箇所は，必ず角の関係とセットで理解する。
\end{itemize}

\subsection*{今後の学習で意識するポイント}
\begin{itemize}
  \item 「接線」「内接」「対角線」などのキーワードから，使う定理や公式を連想できるようにする。
  \item 図形問題では，式を書く前に図に辺や角の情報を書き込んで整理する。
  \item 面積公式，正弦定理，余弦定理の使い分けを練習する。
\end{itemize}

\bigskip

\section*{2次関数}

\subsection*{この分野で起こりやすいミス}
\begin{itemize}
  \item 頂点と最大値・最小値の関係が曖昧である。
  \item 上に凸か下に凸かを取り違える。
  \item 場合分けを図なしで処理して混乱する。
\end{itemize}

\subsection*{例題}
\[
y=2x^2-8x+5
\]
について，$0\leqq x\leqq 3$ における最大値・最小値を求める問題や，  
頂点が
\[
(-1,3)
\]
で，ある点を通ることから
\[
f(x)=a(x+1)^2+3
\]
とおいて係数 $a$ を決定する問題。

\subsection*{例題で起こりがちなミス}
\begin{itemize}
  \item 頂点が定義域の内側にあるか外側にあるかを確認しない。
  \item 最大値・最小値の条件から，$a>0$ か $a<0$ かを判断できない。
  \item いきなり
  \[
  ax^2+bx+c
  \]
  の形で考えてしまい，見通しが悪くなる。
\end{itemize}

\subsection*{解説を読むときに注意してほしいこと}
\begin{itemize}
  \item まず頂点の位置を確認し，その後に定義域との関係を見ている流れを追う。
  \item 頂点形式
  \[
  y=a(x-p)^2+q
  \]
  を採用している理由を理解する。
  \item 「この点を通るから代入する」という係数決定の手順を，型として読む。
\end{itemize}

\subsection*{今後の学習で意識するポイント}
\begin{itemize}
  \item 2次関数は，まず頂点形式で見る習慣をつける。
  \item 問題ごとに必ずグラフを手で描く。
  \item 境界値の前後で場合分けが必要そうなら，グラフを複数描いて考える。
\end{itemize}

\bigskip

\section*{データの分析}

\subsection*{この分野で起こりやすいミス}
\begin{itemize}
  \item 図から読めることまで計算しようとして，かえって遠回りしてしまう。
  \item 相関係数や外れ値の公式を，意味を理解しないまま暗記している。
  \item 分散と四分位範囲を，同じ「ばらつき」として大まかに扱ってしまう。
\end{itemize}

\subsection*{例題}
散布図から
\[
T_{\text{前}}<470,\qquad T_{\text{後}}\geqq 460
\]
を満たす人数を読み取る問題や，相関係数を
\[
r=\frac{\text{共分散}}{\text{標準偏差}\times \text{標準偏差}}
\]
で求める問題。  
また，外れ値の基準
\[
Q_1-1.5(Q_3-Q_1),\qquad Q_3+1.5(Q_3-Q_1)
\]
から四分位範囲を逆算する問題。

\subsection*{例題で起こりがちなミス}
\begin{itemize}
  \item 散布図の点の数を正しく数えられない。
  \item 相関係数で，分母の標準偏差を1つしか掛けない。
  \item 外れ値の式を順方向でしか見られず，$Q_1,Q_3$ を逆算できない。
\end{itemize}

\subsection*{解説を読むときに注意してほしいこと}
\begin{itemize}
  \item まず「図から読み取っているだけなのか」と「計算して求めているのか」を区別する。
  \item 相関係数は，数値そのものよりも式の構造を確認する。
  \item 外れ値の問題では，「結果として与えられた境界値から何を逆算しているのか」を追う。
\end{itemize}

\subsection*{今後の学習で意識するポイント}
\begin{itemize}
  \item 散布図や箱ひげ図では，まず目で読み取れる情報を整理する。
  \item 平均，分散，標準偏差，相関係数の意味を，自分の言葉で説明できるようにする。
  \item 「ばらつき」を表す指標には複数の種類があることを意識する。
\end{itemize}

\bigskip

\section*{図形の性質}

\subsection*{この分野で起こりやすいミス}
\begin{itemize}
  \item 空間図形をそのまま眺めてしまい，必要な平面を取り出せない。
  \item 定理を単発で覚えていて，つながりの中で使えない。
  \item 比の向きや対応を取り違える。
\end{itemize}

\subsection*{例題}
内心 $I$ を用いて角の二等分線の性質を使う問題や，
\[
\angle PED=\angle PID
\]
から4点 $P,E,I,D$ が同一円周上にあると判断し，さらに方べきの定理
\[
AE\cdot AP=AI\cdot AD
\]
を使う問題。  
また，メネラウスの定理を用いて比を求め，空間図形の体積比につなげる問題。

\subsection*{例題で起こりがちなミス}
\begin{itemize}
  \item 「角が等しい」という条件から，共円を疑えない。
  \item メネラウスの定理を，どの三角形に対して使えばよいか分からない。
  \item 体積比を，高さの比ではなく辺の比の感覚で処理してしまう。
\end{itemize}

\subsection*{解説を読むときに注意してほしいこと}
\begin{itemize}
  \item 「どの条件からどの定理へ進んでいるか」という流れを見る。
  \item 共円，方べき，メネラウスがどのように連鎖しているかに注目する。
  \item 空間図形では，どの平面を切り出して2次元で見ているのかを確認する。
\end{itemize}

\subsection*{今後の学習で意識するポイント}
\begin{itemize}
  \item 定理は，「使える条件」とセットで覚える。
  \item 空間図形では，必要な断面図を必ず描く。
  \item 比の問題では，対応する線分を図に明記する。
\end{itemize}

\bigskip

\section*{場合の数・確率}

\subsection*{この分野で起こりやすいミス}
\begin{itemize}
  \item 場合分けに漏れや重複が出る。
  \item 掛ける確率と足す確率を混同する。
  \item 先に計算を始めてしまい，全体の構造を見失う。
\end{itemize}

\subsection*{例題}
3人または4人のリーグ戦で，Aが優勝する確率を求める問題。  
たとえば，
\[
\text{「Aが2勝0敗で優勝する場合」}
\]
と
\[
\text{「Aが1勝1敗で，抽選により優勝する場合」}
\]
を分けて考える問題や，4人の場合に
\[
\text{「全敗する人がいる場合」},\qquad
\text{「全敗する人がいない場合」}
\]
で場合分けする問題。

\subsection*{例題で起こりがちなミス}
\begin{itemize}
  \item 抽選で選ばれる確率 $\frac13,\ \frac12$ などを掛け忘れる。
  \item 場合分けの設計が不十分で，特定のパターンを数え漏らす。
  \item 対戦結果の確率と，その後の抽選確率を混同する。
\end{itemize}

\subsection*{解説を読むときに注意してほしいこと}
\begin{itemize}
  \item まず「どのように場合分けしているか」を確認する。
  \item 各場合で，「対戦結果の確率」と「抽選で選ばれる確率」を分けて読む。
  \item 最後に足し合わせているのが，互いに重ならない場合だけであることを意識する。
\end{itemize}

\subsection*{今後の学習で意識するポイント}
\begin{itemize}
  \item いきなり計算せず，最初に場合分けを日本語で書く。
  \item 樹形図や表を使って，全パターンを整理する。
  \item 確率は「構造を作ること」が先で，計算は後だと意識する。
\end{itemize}

\bigskip

\section*{全体のまとめ}

\subsection*{解説を読むうえで大切なこと}
解説は答えを確認するためだけに読むものではなく，
\[
\text{「なぜその発想になったのか」}
\]
を学ぶために読むものである。

\subsection*{特に見るべき3点}
\begin{enumerate}
  \item 条件をどのように言い換えているか
  \item なぜその式や定理を使ったのか
  \item どの順番で処理しているか
\end{enumerate}

\subsection*{今後の学習で共通して意識するポイント}
\begin{itemize}
  \item 計算だけでなく，発想の出発点を言語化する。
  \item 図や表を省略しない。
  \item 「この問題だけできるか」ではなく，「同じ分野でまた出たときに再現できるか」を基準に復習する。
\end{itemize}
